\documentclass[11pt]{article}
\usepackage[margin=1in]{geometry}
\usepackage{amsmath,amssymb,microtype}
\usepackage[colorlinks=true,urlcolor=blue,citecolor=blue]{hyperref}
\newcommand{\pd}{p_{d_1,d_2}}
\title{The Grushin topological-dimension conjecture of arXiv:1406.4025}
\author{Literature-status note for \texttt{fa\_banach\_001}}
\date{26 August 2026}

\begin{document}
\maketitle

\noindent\textbf{Status.}
\texttt{literature\_already\_answered (full affirmative)}.

\section*{Original conjecture}

Chen and Ouhabaz study the Grushin operator
\[
 L=-\Delta_x-|x|^2\Delta_y
 \quad\text{on }\mathbb R^{d_1}\times\mathbb R^{d_2}.
\]
Their Theorems 1.1 and 1.2 use
\[
 D=\max\{d_1+d_2,2d_2\}
 \quad\text{and}\quad
 \pd=\min\left\{\frac{2d_1}{d_1+2},
                    \frac{2(d_2+1)}{d_2+3}\right\}.
\]
In the unnumbered Conjecture on page 3 of the arXiv PDF, immediately after
those theorems, they write that the previous theorems should hold with
$D=d_1+d_2$ instead of $D=\max\{d_1+d_2,2d_2\}$ \cite{ChenOuhabaz}.
Thus the open part asks for both:
\begin{enumerate}
\item the $p$-specific spectral multiplier threshold
  $s>(d_1+d_2)(1/p-1/2)$ in the source range; and
\item uniform $L^p$ boundedness of $(1-tL)_+^\delta$ when
  $\delta>(d_1+d_2)(1/p-1/2)-1/2$ in the same range.
\end{enumerate}
The source observes that the multiplier assertion was already known for
$p=1$; the unresolved regime was the sharp $p$-specific improvement for
$p>1$ when $d_2>d_1$.

\section*{Answer}

Niedorf explicitly identifies the expectation of Chen and Ouhabaz and cites
their paper before stating his main results \cite{Niedorf}. On page 3 of the
arXiv PDF, Theorem 1.1 proves that for $1<p\leq\pd$, every bounded Borel
function with local Sobolev regularity
\[
 s>(d_1+d_2)(1/p-1/2)
\]
defines a bounded spectral multiplier $F(L)$ on $L^p$. Theorem 1.2 on the
same page proves that for $1\leq p\leq\pd$ and
\[
 \delta>(d_1+d_2)(1/p-1/2)-1/2,
\]
the Bochner--Riesz means $(1-tL)_+^\delta$ are uniformly bounded on $L^p$.
These are exactly the two source theorems with $D$ replaced by the
topological dimension $d_1+d_2$. Together with the already-known $p=1$
multiplier endpoint recorded by the source, this is a full affirmative answer
to the displayed conjecture.

The answering authors knew the relation: the abstract says their result
improves the Chen--Ouhabaz theorem, and the introduction expressly attributes
the expected replacement to Chen and Ouhabaz before stating Theorems 1.1 and
1.2.

\section*{Scope and search evidence}

This packet concerns only the stated Chen--Ouhabaz range
$1\leq p\leq\pd$ and the two conclusions in their unnumbered Conjecture. It
does not assert a wider $p$-range or resolve the classical Euclidean
Bochner--Riesz conjecture.

The four run indexes were searched for arXiv:1406.4025 and the terms
``Grushin'', ``spectral multiplier'', and ``Bochner--Riesz''; they contained
related Grushin packets but no record for this exact conjecture. Exact arXiv
and title/keyword searches located arXiv:2110.10058. The identification was
then checked in the two PDFs: the source Conjecture is on page 3, and the
answering discussion and Theorems 1.1-1.2 are on page 3 of the supporting
paper. No new mathematical result is claimed.

\begin{thebibliography}{9}
\bibitem{ChenOuhabaz}
P.~Chen and E.~M. Ouhabaz,
\emph{Weighted restriction type estimates for Grushin operators and
application to spectral multipliers and Bochner--Riesz summability},
Math. Z. \textbf{282} (2016), 663--678;
arXiv:1406.4025,
\url{https://arxiv.org/abs/1406.4025}.

\bibitem{Niedorf}
L.~Niedorf,
\emph{A $p$-specific spectral multiplier theorem with sharp regularity bound
for Grushin operators},
Math. Z. \textbf{301} (2022), 4153--4173;
arXiv:2110.10058,
\url{https://arxiv.org/abs/2110.10058}.
\end{thebibliography}

\end{document}

